\documentclass[12pt,a4paper]{article}
\usepackage[margin=2.5cm]{geometry}
\usepackage{amsmath,amssymb,amsthm}
\usepackage{setspace}
\usepackage{microtype}
\usepackage{xcolor}
\usepackage{hyperref}
\hypersetup{colorlinks=true,linkcolor=blue!50!black,urlcolor=blue!50!black,citecolor=blue!50!black}
\onehalfspacing

\newtheorem{theorem}{Theorem}
\newtheorem{lemma}[theorem]{Lemma}
\newtheorem{proposition}[theorem]{Proposition}
\newtheorem*{remark}{Remark}

\title{\textbf{On the Unique Catalan Identity in Ramanujan's Nested Radical Family}}
\author{Yan Senez\\\texttt{yan.senez@protonmail.com}}
\date{\today}

\begin{document}
\maketitle

\begin{abstract}
For each integer $n \geq 1$, Ramanujan's nested radical identity
\[
n + 1 \;=\; \sqrt{\,1 + n\sqrt{1 + (n+1)\sqrt{1 + (n+2)\sqrt{\cdots}}}\,}
\]
holds, with first-level unfolding $(n+1)^2 = 1 + n(n+2)$.  We show
that the case $n = 2$, yielding $3^2 - 2^3 = 1$, is the unique member
of the family for which the quantity $n(n+2)$ is a non-trivial perfect
power.  The proof reduces to two elementary observations and an
application of Mihailescu's theorem~\cite{Mihailescu2004}, the
resolution of Catalan's conjecture.  We comment on what this says
about the place of the case $n=2$ within the Ramanujan family.
\end{abstract}

\section*{1. Introduction}

In a problem submitted to the \emph{Journal of the Indian Mathematical
Society}~\cite{Ramanujan1911}, Ramanujan asked the reader to evaluate
\begin{equation}
\sqrt{\,1 + 2\sqrt{1 + 3\sqrt{1 + 4\sqrt{1 + 5\sqrt{1 + \cdots}}}}\,},
\label{eq:Rama-3}
\end{equation}
giving the answer $3$.  This nested radical is one of the most famous
identities in elementary arithmetic.  It is in fact a special case of
the general formula recorded in Ramanujan's
\emph{Notebooks}~\cite[Ch.~12]{Berndt1985}:
\begin{equation}
n + 1 \;=\; \sqrt{\,1 + n\sqrt{1 + (n+1)\sqrt{1 + (n+2)\sqrt{\cdots}}}\,},
\qquad n \in \mathbb{N}_{\geq 1}.
\label{eq:Rama-general}
\end{equation}
The convergence of the right-hand side as a limit of finite truncations
was settled by Vijayaraghavan~\cite{Vijayaraghavan1929} and
Herschfeld~\cite{Herschfeld1935}; the algebraic identity itself follows
from the recursive relation $f(n)^2 = 1 + n \cdot f(n+1)$ with $f(n) = n+1$,
since $(n+1)^2 - 1 = n(n+2)$.

The case $n = 2$ in (\ref{eq:Rama-general}) recovers (\ref{eq:Rama-3}).
Setting aside its visual elegance, it is natural to ask whether $n = 2$
enjoys any \emph{intrinsic} singularity within the family
(\ref{eq:Rama-general}).  At first glance, no: the identity is
parametric, and every $n \geq 1$ produces a perfectly valid nested
radical.  In this note we observe that, when one looks at the
first-level expansion
\begin{equation}
(n+1)^2 \;=\; 1 + n(n+2),
\label{eq:first-level}
\end{equation}
the case $n = 2$ does enjoy a singular property: it is the
\emph{unique} value of $n$ for which $n(n+2)$ is a non-trivial perfect
power, and accordingly the unique value for which (\ref{eq:first-level})
is itself a Catalan-type identity
\begin{equation}
3^2 - 2^3 \;=\; 1.
\label{eq:Catalan-3-2}
\end{equation}

The core of the observation is that this uniqueness is not merely
empirical: it is a consequence of Mihailescu's
theorem~\cite{Mihailescu2004}, the proof of Catalan's conjecture.  The
purpose of this note is to record this remark, give a clean proof, and
discuss what it says about the place of $n=2$ in the Ramanujan family.

\section*{2. The Result}

\begin{theorem}\label{thm:main}
For $n \in \mathbb{Z}_{\geq 1}$, the integer $n(n+2)$ is a non-trivial
perfect power, i.e.\ there exist $b, q \in \mathbb{Z}_{\geq 2}$ with
\[
n(n+2) \;=\; b^{q},
\]
if and only if $n = 2$, in which case $(b, q) = (2, 3)$ and the
identity~(\ref{eq:first-level}) reduces to
$3^{2} - 2^{3} = 1$.
\end{theorem}

\begin{proof}
We split according to the value of $q \geq 2$.

\medskip
\noindent\textbf{Case 1: $q = 2$.}  Suppose $n(n+2) = b^2$ for some
$b \geq 2$.  Using $n(n+2) = (n+1)^2 - 1$, this is
\begin{equation}
(n+1)^2 - b^2 \;=\; 1,
\quad\text{i.e.,}\quad
\bigl((n+1) - b\bigr)\bigl((n+1) + b\bigr) \;=\; 1.
\label{eq:diff-squares}
\end{equation}
Since $b \geq 2$ and $n \geq 1$, both factors on the left of
(\ref{eq:diff-squares}) are positive integers; the only factorisation
of $1$ in $\mathbb{Z}_{>0}$ is $1 \cdot 1$, forcing
$(n+1) - b = (n+1) + b = 1$, hence $b = 0$ and $n + 1 = 1$,
contradicting $b \geq 2$ and $n \geq 1$.  Thus no $n \geq 1$ makes
$n(n+2)$ a perfect square.

(Equivalently: $(n+1)^2 - 1$ lies strictly between the consecutive
squares $n^2$ and $(n+1)^2$ for $n \geq 1$, hence is not a square.)

\medskip
\noindent\textbf{Case 2: $q \geq 3$.}  Suppose $n(n+2) = b^q$ for some
$b \geq 2$ and $q \geq 3$.  Then
\begin{equation}
(n+1)^2 - b^q \;=\; 1,
\label{eq:cat-shape}
\end{equation}
which is the Catalan equation $a^p - b^q = 1$ with parameters
$(a, p, q) = (n+1, 2, q)$.  By Mihailescu's theorem
\cite{Mihailescu2004}, the unique solution in positive integers to
\[
a^p - b^q = 1, \qquad a, b, p, q \geq 2,
\]
is $(a, b, p, q) = (3, 2, 2, 3)$.  Specialising to our case
($p = 2$, $q \geq 3$) yields $a = n+1 = 3$, $b = 2$, $q = 3$, and
therefore $n = 2$ with $n(n+2) = 2 \cdot 4 = 8 = 2^3$.

\medskip
Combining the two cases, $n(n+2) = b^q$ with $b, q \geq 2$ forces
$n = 2$, $b = 2$, $q = 3$, and (\ref{eq:first-level}) becomes the
identity $3^2 - 2^3 = 1$.  Conversely, $n = 2$ obviously gives such a
representation.
\end{proof}

\section*{3. Remarks}

\subsection*{3.1.\ Depth of the theorem versus simplicity of the application.}

Theorem~\ref{thm:main} reduces to a one-line application of
Mihailescu~\cite{Mihailescu2004}, but the underlying theorem is deep:
its proof rests on the theory of cyclotomic units and primary roots
in cyclotomic fields, building on earlier reductions by Cassels,
Inkeri and others.  What is novel here is not the depth, but the
observation that this depth manifests itself in an entirely elementary
context, that of Ramanujan's nested radicals.  The theorem is, in a
sense, a translation of Mihailescu into the Ramanujan family.

\subsection*{3.2.\ Higher levels of the radical.}

The first-level expansion (\ref{eq:first-level}) is special.  At depth
$k$, repeated unfolding gives
\[
(n+1)^2 \;=\; 1 + n \, f(n+1)
\;=\; 1 + n\sqrt{1 + (n+1)\sqrt{1 + (n+2)\sqrt{\cdots}}},
\]
and the values appearing inside no longer have the structure of
$(n+1)^2 - 1$.  Whether some \emph{deeper} level of the radical for
some other $n$ produces a Catalan-type identity is a separate
question; we do not address it here, but note that it would require a
non-elementary perfect-power identity in $\sqrt{1 + (n+1)\sqrt{\cdots}}$,
which is not known to occur.

\subsection*{3.3.\ Connection to Pillai's conjecture.}

Pillai's conjecture asserts that for every $k \geq 1$, the equation
$a^p - b^q = k$ has finitely many solutions in $(a, b, p, q)$ with
$a, b, p, q \geq 2$.  The Catalan case $k = 1$ is now known
(Mihailescu); higher $k$ remain open.  Our observation can be
rephrased as: \emph{within the Ramanujan family}~(\ref{eq:Rama-general}),
the perfect-power difference $|a^p - b^q|$ at the first level reaches
the value $1$ exactly once, at $n = 2$.  Whether the same can be said
for higher gap values $|a^p - b^q| = k$, $k \geq 2$, depends on
parametric solutions of $n(n+2) - b^q = \pm k$, where partial results
(e.g.\ Bilu, Bugeaud) apply but no general analogue of Mihailescu's
theorem is available.

\subsection*{3.4.\ Visual interpretation.}

Theorem~\ref{thm:main} can be paraphrased as follows.  The family
(\ref{eq:Rama-general}) realises every positive integer $n+1$ as a
nested radical; among these, the integer $3$ is the only one whose
\emph{innermost defining product} ($2 \cdot 4$) collapses
multiplicatively into a single perfect power ($2^3$).  This is what
makes the case $n=2$, of all the visually similar cases in the family,
arithmetically singular.

\begin{thebibliography}{9}

\bibitem{Berndt1985}
B.\ C.\ Berndt,
\emph{Ramanujan's Notebooks, Part~II},
Springer-Verlag, New York, 1989.

\bibitem{Herschfeld1935}
A.\ Herschfeld,
\emph{On infinite radicals},
Amer.\ Math.\ Monthly \textbf{42} (1935), no.\ 7, 419--429.

\bibitem{Mihailescu2004}
P.\ Mihailescu,
\emph{Primary cyclotomic units and a proof of Catalan's conjecture},
J.\ reine angew.\ Math.\ \textbf{572} (2004), 167--195.

\bibitem{Ramanujan1911}
S.\ Ramanujan,
\emph{Question 289},
J.\ Indian Math.\ Soc.\ \textbf{3} (1911), 90;
\emph{Solution to Question 289},
J.\ Indian Math.\ Soc.\ \textbf{4} (1912), 226.

\bibitem{Vijayaraghavan1929}
T.\ Vijayaraghavan,
\emph{On certain infinite series and convergent integrals},
Proc.\ London Math.\ Soc.\ (2) \textbf{29} (1929), 281--294.

\end{thebibliography}

\end{document}
